% Purpose: PhD Dissertation Abstract summarizing research
% motivations, methods, benchmark datasets, HPC implementation, and
% scientific outcomes.
% Status: REVISED
% Source-of-truth: OGS/src/, doc/MHPC/chapters/abstract.tex,
% experimental results.

\begin{abstract}
  The exponential increase of semi-continuous seismic recording archives and
  the deployment of dense regional networks have created severe operational
  bottlenecks for traditional earthquake monitoring workflows. Classical
  automated phase pickers, such as \ac{STA/LTA} and autoregressive \ac{AIC}
  formulations \eg{\parencite{TakanamiKitagawa1988}}, degrade substantially
  when confronted with low \ac{SNR}, emergent phase onsets, overlapping seismic
  codas, and anthropogenic noise. Consequently, regional seismic observatories
  remain heavily dependent on labor-intensive manual analyst review,
  introducing latency and limiting catalog completeness at microseismic scales.

  This dissertation presents the design, implementation, and empirical
  validation of an end-to-end, procedural \ac{ML} pipeline for automated
  earthquake pick detection, phase association, hypocentral relocation, and
  magnitude estimation, specifically engineered for scalable execution on
  modern \ac{HPC} architectures. The computational framework leverages the
  open-source SeisBench \parencite{SeisBench} ecosystem to deploy \ac{DL} phase
  picking architectures, including the 1D U-Net convolutional network
  \ac{PhaseNet} and the hierarchical convolutional-transformer network
  \ac{EQT}. To systematically investigate \textbf{domain adaptation} and
  \textbf{model transferability}, these pickers are evaluated across diverse
  regional and global benchmark training datasets: the Italian seismic dataset
  \ac{INSTANCE}; the global seismic dataset \ac{STEAD}; the Southern California
  seismic dataset \ac{SCEDC}; and the original author training sets.

  The target case study comprises continuous three-component waveform streams
  recorded by the \acf{SMINO}, managed by the \acf{OGS}, augmented by seventeen
  regional and transboundary networks in the surrounding regions. The network
  configuration encompasses \placeholder{NUM STATIONS} stations. The
  observation horizon spans over twenty years of observational data
  (\textbf{January 1st, 2005 to December 31st, 2025}), capturing multi-decade
  regional background microseismicity as well as prominent seismic sequences
  including the March 27, 2024, $M_{\mathrm{L}}~\num{4.6}$
  ($M_{\mathrm{w}}~\num{4.5}$) Socchieve earthquake in Carnia, Friuli, one of
  the most energetic seismic episodes in Northeastern Italy since the
  destructive 1976 sequence.

  Detected phase arrivals are associated into discrete seismic events using two
  complementary state-of-the-art algorithms: the Bayesian \ac{GaMMA} and the
  high-throughput octree-based associator \ac{PyOcto}. Associated events are
  relocated using a 1D probabilistic non-linear search via \ac{NonLinLoc},
  followed by a parametric local magnitude ($M_{\mathrm{L}}$) estimation. To
  establish an unbiased comparative baseline, the \ac{OGS} reference catalog
  for the observation window was independently relocated under the identical 1D
  \ac{NonLinLoc} velocity model, yielding \placeholder{NUM EVENTS} and
  \placeholder{NUM PICKS} manually curated phase picks
  (\placeholder{NUM P PICKS}~$P$ and \placeholder{NUM S PICKS}~$S$).
  Performance is evaluated using a physically bounded \ac{BPGMA}, resolving
  matched, missed, and proposed events and picks with strict spatiotemporal and
  phase-type conservation.

  Distributed execution on the \ac{CINECA} Leonardo supercomputer (EuroHPC
    Booster partition, utilizing NVIDIA Ampere A100 \acp{GPU} and multi-core
  \ac{CPU} pipelines) demonstrates near-linear scalability, reducing continuous
  waveform processing times by orders of magnitude compared to legacy serial
  workflows. Systematic evaluation reveals that \ac{PhaseNet} trained on the
  regional \ac{INSTANCE} dataset achieves the highest pick recovery, optimal
  phase residual distributions, and the most reliable association completeness.
  The proposed infrastructure provides operational observatories with a
  scalable, reproducible foundation for high-resolution historical catalog
  reprocessing.
\end{abstract}
